% =============================================================
% Matriz de Rastreabilidade de Testes — Sprint 1
% GreenHerb API  |  Engenharia de Software II
% Compilar com: pdflatex matriz_rastreabilidade.tex
% =============================================================
\documentclass[a4paper,11pt]{article}

\usepackage[utf8]{inputenc}
\usepackage[T1]{fontenc}
\usepackage[portuguese]{babel}
\usepackage[a4paper, top=2.5cm, bottom=2.5cm, left=2.5cm, right=2.5cm]{geometry}
\usepackage{booktabs}
\usepackage{longtable}
\usepackage{array}
\usepackage[table, dvipsnames]{xcolor}
\usepackage{pdflscape}
\usepackage{hyperref}
\usepackage{parskip}
\usepackage{titlesec}
\usepackage{fancyhdr}
\usepackage{microtype}

% ------------------------------------------------------------------
% Cores
% ------------------------------------------------------------------
\definecolor{headerblue}{RGB}{68,114,196}
\definecolor{rowalt}{RGB}{235,241,250}
\definecolor{grouprow}{RGB}{217,225,242}
\definecolor{titlecolor}{RGB}{31,73,125}

% ------------------------------------------------------------------
% Formatação de secções
% ------------------------------------------------------------------
\titleformat{\section}{\large\bfseries\color{titlecolor}}{}{0em}{}[\color{headerblue}\titlerule]
\titleformat{\subsection}{\normalsize\bfseries\color{titlecolor}}{}{0em}{}

% ------------------------------------------------------------------
% Novos tipos de coluna (já com \footnotesize)
% ------------------------------------------------------------------
\newcolumntype{L}[1]{>{\raggedright\arraybackslash\footnotesize}p{#1}}
\newcolumntype{C}[1]{>{\centering\arraybackslash\footnotesize}p{#1}}

% ------------------------------------------------------------------
% Espaçamento das células
% ------------------------------------------------------------------
\setlength{\tabcolsep}{3pt}
\renewcommand{\arraystretch}{1.35}

% ------------------------------------------------------------------
% Cabeçalho de página
% ------------------------------------------------------------------
\pagestyle{fancy}
\fancyhf{}
\rhead{\footnotesize GreenHerb API — Matriz de Rastreabilidade}
\lhead{\footnotesize Sprint 1 — Módulo de Autenticação}
\cfoot{\footnotesize \thepage}
\renewcommand{\headrulewidth}{0.4pt}

\hypersetup{
  colorlinks = true,
  linkcolor  = headerblue,
  urlcolor   = headerblue,
  pdftitle   = {Matriz de Rastreabilidade — Sprint 1},
}

% =============================================================
\begin{document}
% =============================================================

% ------------------------------------------------------------------
% Título
% ------------------------------------------------------------------
\begin{center}
  {\LARGE\bfseries\color{titlecolor} GreenHerb API}\\[6pt]
  {\large\bfseries Matriz de Rastreabilidade de Testes}\\[4pt]
  {\normalsize Sprint 1 — Módulo de Autenticação}\\[10pt]
  \begin{tabular}{@{}ll@{}}
    \textbf{UC:}    & Engenharia de Software II — Testes de Software \\[2pt]
    \textbf{Data:}  & 2026-05-07 \\
  \end{tabular}
\end{center}

\vspace{0.5cm}
\tableofcontents
\newpage

% =============================================================
\section{Legenda de Prefixos}
% =============================================================

\begin{tabular}{@{}C{2cm}L{9cm}@{}}
  \toprule
  \textbf{Prefixo} & \textbf{Nível de Teste} \\
  \midrule
  \rowcolor{rowalt} TU & Teste de Unidade \\
  TI & Teste de Integração \\
  \rowcolor{rowalt} TS & Teste de Sistema \\
  \bottomrule
\end{tabular}

% =============================================================
\section{Requisitos Funcionais de Referência}
% =============================================================

\begin{tabular}{@{}C{1.4cm}L{13cm}@{}}
  \toprule
  \textbf{ID} & \textbf{Descrição} \\
  \midrule
  \rowcolor{rowalt}
  RF-01 & O sistema deve autenticar utilizadores com \texttt{username} e \texttt{password},
          devolvendo tokens JWT. \\
  RF-02 & O sistema deve gerar um \textit{access token} JWT com payload \texttt{\{id, perfil\}}
          e expiração de 2\,h. \\
  \rowcolor{rowalt}
  RF-03 & O sistema deve gerar um \textit{refresh token} JWT com payload \texttt{\{id, perfil\}}
          e expiração de 7\,d. \\
  RF-04 & O sistema deve renovar o \textit{access token} a partir de um \textit{refresh token} válido. \\
  \rowcolor{rowalt}
  RF-05 & O sistema deve verificar a validade e integridade de um token JWT. \\
  RF-06 & O sistema deve controlar o acesso por perfil
          (\texttt{Técnico}, \texttt{Responsável}, \texttt{Administrador}). \\
  \bottomrule
\end{tabular}

% =============================================================
\section{Itens de Teste}
% =============================================================

\begin{tabular}{@{}C{1.2cm}L{4cm}C{2.5cm}L{3cm}L{5cm}@{}}
  \toprule
  \textbf{ID} & \textbf{Função} & \textbf{Req.} & \textbf{Técnicas} & \textbf{Ficheiro} \\
  \midrule
  \rowcolor{rowalt}
  TI-01 & \texttt{generateAccessToken}  & RF-02            & PE + CM & \texttt{src/services/authService.js} \\
  TI-02 & \texttt{generateRefreshToken} & RF-03            & PE + CM & \texttt{src/services/authService.js} \\
  \rowcolor{rowalt}
  TI-03 & \texttt{verifyToken}          & RF-05            & PE      & \texttt{src/services/authService.js} \\
  TI-04 & \texttt{renewAccessToken}     & RF-04            & PE      & \texttt{src/services/authService.js} \\
  \rowcolor{rowalt}
  TI-05 & \texttt{login}                & RF-01, RF-02, RF-03 & PE + CM & \texttt{src/services/authService.js} \\
  TI-06 & \texttt{hasProfile}           & RF-06            & PE + CM & \texttt{src/services/authService.js} \\
  \bottomrule
\end{tabular}

% =============================================================
\section{Matriz de Rastreabilidade}
% =============================================================

% Larguras calculadas para página A4 em landscape com margens de 2,5 cm:
%   largura disponível ≈ 247 mm = 24,7 cm
%   soma das colunas: 0,9+1,1+3,2+1,4+3,8+5,8+5,8 = 22,0 cm
%   espaço de separação (6 gaps × 3pt × 2 = 36pt ≈ 1,27 cm)
%   total ≈ 23,27 cm < 24,7 cm  ✓

\begin{landscape}
\setlength{\tabcolsep}{3pt}
\renewcommand{\arraystretch}{1.4}

\begin{longtable}{
  C{0.9cm}   % ID CT
  C{1.1cm}   % Requisito
  L{3.2cm}   % Endpoint(s)
  C{1.4cm}   % Nível
  L{3.8cm}   % Técnica
  L{5.8cm}   % Resultado Esperado
  L{5.8cm}   % Pré-condições
}

\caption{Casos de Teste — Sprint~1 (Módulo de Autenticação)}
\label{tab:matriz}\\

% Cabeçalho — primeira página
\toprule
\textbf{\footnotesize ID CT} &
\textbf{\footnotesize Req.} &
\textbf{\footnotesize Endpoint(s) Exercitado(s)} &
\textbf{\footnotesize Nível} &
\textbf{\footnotesize Técnica Aplicada} &
\textbf{\footnotesize Resultado Esperado} &
\textbf{\footnotesize Pré-condições} \\
\midrule
\endfirsthead

% Cabeçalho — páginas seguintes
\multicolumn{7}{l}{\footnotesize\textit{(continuação da Tabela~\ref{tab:matriz})}}\\[3pt]
\toprule
\textbf{\footnotesize ID CT} &
\textbf{\footnotesize Req.} &
\textbf{\footnotesize Endpoint(s) Exercitado(s)} &
\textbf{\footnotesize Nível} &
\textbf{\footnotesize Técnica Aplicada} &
\textbf{\footnotesize Resultado Esperado} &
\textbf{\footnotesize Pré-condições} \\
\midrule
\endhead

% Rodapé intermédio
\midrule
\multicolumn{7}{r}{\footnotesize\textit{continua na página seguinte}}\\
\endfoot

% Rodapé final
\bottomrule
\endlastfoot

% ==============================================================
% TI-01  —  generateAccessToken
% ==============================================================
\rowcolor{grouprow}
\multicolumn{7}{l}{%
  \footnotesize\textbf{TI-01 — \texttt{generateAccessToken(userId, perfil)}}
  \quad Condição analisada: \texttt{!userId || !perfil}}\\
\midrule

\rowcolor{white}
TU-01 & RF-02 &
\texttt{POST /auth/login} (indirecto) &
Unidade &
CM\,—\,cond.\,(T$|$T) &
Lança \texttt{Error}: \textit{"userId and perfil are required"} &
\texttt{JWT\_SECRET} definido em \texttt{process.env}; sem dependência de BD \\

\rowcolor{rowalt}
TU-02 & RF-02 &
\texttt{POST /auth/login} (indirecto) &
Unidade &
CM\,—\,cond.\,(T$|$F) &
Lança \texttt{Error}: \textit{"userId and perfil are required"} &
Idem \\

\rowcolor{white}
TU-03 & RF-02 &
\texttt{POST /auth/login} (indirecto) &
Unidade &
CM\,—\,cond.\,(F$|$T) &
Lança \texttt{Error}: \textit{"userId and perfil are required"} &
Idem \\

\rowcolor{rowalt}
TU-04 & RF-02 &
\texttt{POST /auth/login} (indirecto) &
Unidade &
CM (F$|$F) + PE classe válida &
Retorna string JWT; payload verificável contém \texttt{id} e \texttt{perfil} correctos &
\texttt{JWT\_SECRET} definido; \texttt{userId} e \texttt{perfil} não nulos \\

\midrule
% ==============================================================
% TI-02  —  generateRefreshToken
% ==============================================================
\rowcolor{grouprow}
\multicolumn{7}{l}{%
  \footnotesize\textbf{TI-02 — \texttt{generateRefreshToken(userId, perfil)}}
  \quad Condição analisada: \texttt{!userId || !perfil}}\\
\midrule

\rowcolor{white}
TU-05 & RF-03 &
\texttt{POST /auth/login} (indirecto) &
Unidade &
CM\,—\,cond.\,(T$|$T) &
Lança \texttt{Error}: \textit{"userId and perfil are required"} &
\texttt{JWT\_REFRESH\_SECRET} definido em \texttt{process.env} \\

\rowcolor{rowalt}
TU-06 & RF-03 &
\texttt{POST /auth/login} (indirecto) &
Unidade &
CM\,—\,cond.\,(T$|$F) &
Lança \texttt{Error}: \textit{"userId and perfil are required"} &
Idem \\

\rowcolor{white}
TU-07 & RF-03 &
\texttt{POST /auth/login} (indirecto) &
Unidade &
CM\,—\,cond.\,(F$|$T) &
Lança \texttt{Error}: \textit{"userId and perfil are required"} &
Idem \\

\rowcolor{rowalt}
TU-08 & RF-03 &
\texttt{POST /auth/login} (indirecto) &
Unidade &
CM (F$|$F) + PE classe válida &
Retorna JWT assinado com \texttt{JWT\_REFRESH\_SECRET}; payload contém \texttt{id} e \texttt{perfil} &
\texttt{JWT\_REFRESH\_SECRET} e \texttt{userId}/\texttt{perfil} não nulos \\

\rowcolor{white}
TU-09 & RF-03 &
\texttt{POST /auth/login} (indirecto) &
Unidade &
PE — fallback de \textit{secret} &
Retorna JWT assinado com \texttt{JWT\_SECRET} quando \texttt{JWT\_REFRESH\_SECRET} não está definido &
\texttt{JWT\_REFRESH\_SECRET} removido temporariamente do ambiente \\

\midrule
% ==============================================================
% TI-03  —  verifyToken
% ==============================================================
\rowcolor{grouprow}
\multicolumn{7}{l}{%
  \footnotesize\textbf{TI-03 — \texttt{verifyToken(token, secret?)}}
  \quad Condição analisada: \texttt{!token}}\\
\midrule

\rowcolor{white}
TU-10 & RF-05 &
\texttt{POST /auth/refresh}; rotas protegidas &
Unidade &
PE — classe inválida (token nulo) &
Lança \texttt{Error}: \textit{"Token is required"} &
Nenhuma \\

\rowcolor{rowalt}
TU-11 & RF-05 &
\texttt{POST /auth/refresh}; rotas protegidas &
Unidade &
PE — classe inválida (token adulterado) &
Lança \texttt{JsonWebTokenError} &
\texttt{JWT\_SECRET} definido \\

\rowcolor{white}
TU-12 & RF-05 &
\texttt{POST /auth/refresh}; rotas protegidas &
Unidade &
PE — classe inválida (token expirado) &
Lança \texttt{TokenExpiredError} &
Token gerado com \texttt{expiresIn:\,-1} (já expirado na criação) \\

\rowcolor{rowalt}
TU-13 & RF-05 &
\texttt{POST /auth/refresh}; rotas protegidas &
Unidade &
PE — classe válida &
Retorna payload com \texttt{id} e \texttt{perfil} correctos &
Token válido assinado com \texttt{JWT\_SECRET} \\

\rowcolor{white}
TU-14 & RF-05 &
\texttt{POST /auth/refresh}; rotas protegidas &
Unidade &
PE — fallback de \textit{secret} &
Usa \texttt{JWT\_SECRET} do ambiente quando argumento \texttt{secret} é omitido &
Token assinado com \texttt{JWT\_SECRET}; chamada sem argumento \texttt{secret} \\

\midrule
% ==============================================================
% TI-04  —  renewAccessToken
% ==============================================================
\rowcolor{grouprow}
\multicolumn{7}{l}{%
  \footnotesize\textbf{TI-04 — \texttt{renewAccessToken(refreshToken)}}}\\
\midrule

\rowcolor{white}
TU-15 & RF-04 &
\texttt{POST /auth/refresh} &
Unidade &
PE — classe válida &
Retorna novo \textit{access token}; payload verificável contém \texttt{id} e \texttt{perfil} do \textit{refresh token} original &
\textit{Refresh token} válido assinado com \texttt{JWT\_REFRESH\_SECRET} \\

\rowcolor{rowalt}
TU-16 & RF-04 &
\texttt{POST /auth/refresh} &
Unidade &
PE — classe inválida &
Propaga erro de \texttt{verifyToken} (token malformado) &
Token inválido passado como argumento \\

\rowcolor{white}
TU-17 & RF-04 &
\texttt{POST /auth/refresh} &
Unidade &
PE — fallback de \textit{secret} &
Usa \texttt{JWT\_SECRET} quando \texttt{JWT\_REFRESH\_SECRET} ausente; novo \textit{access token} verificável &
\texttt{JWT\_REFRESH\_SECRET} removido temporariamente; \textit{refresh token} assinado com \texttt{JWT\_SECRET} \\

\midrule
% ==============================================================
% TI-05  —  login
% ==============================================================
\rowcolor{grouprow}
\multicolumn{7}{l}{%
  \footnotesize\textbf{TI-05 — \texttt{login(username, password, findUser)}}
  \quad Condições analisadas: (1)~\texttt{!username || !password} \quad (2)~\texttt{!user} \quad (3)~\texttt{!valid}}\\
\midrule

\rowcolor{white}
TU-18 & RF-01 &
\texttt{POST /auth/login} &
Unidade &
CM\,—\,cond.\,1\,(T$|$T) &
Lança \texttt{Error}: \textit{"Username and password are required"} &
\texttt{JWT\_SECRET} definido; \texttt{findUser} não invocado \\

\rowcolor{rowalt}
TU-19 & RF-01 &
\texttt{POST /auth/login} &
Unidade &
CM\,—\,cond.\,1\,(T$|$F) &
Lança \texttt{Error}: \textit{"Username and password are required"} &
Idem \\

\rowcolor{white}
TU-20 & RF-01 &
\texttt{POST /auth/login} &
Unidade &
CM\,—\,cond.\,1\,(F$|$T) &
Lança \texttt{Error}: \textit{"Username and password are required"} &
Idem \\

\rowcolor{rowalt}
TU-21 & RF-01 &
\texttt{POST /auth/login} &
Unidade &
PE — cond.\,2 (utilizador inexistente) &
Lança \texttt{Error}: \textit{"Invalid credentials"} &
\texttt{findUser} mockado a devolver \texttt{null} \\

\rowcolor{white}
TU-22 & RF-01 &
\texttt{POST /auth/login} &
Unidade &
PE — cond.\,3 (password incorrecta) &
Lança \texttt{Error}: \textit{"Invalid credentials"} &
\texttt{findUser} mockado com utilizador válido; \textit{hash} gerado em \texttt{beforeAll} com \textit{bcryptjs} (salt\,=\,10) \\

\rowcolor{rowalt}
TU-23 & RF-01, RF-02, RF-03 &
\texttt{POST /auth/login} &
Unidade &
CM (F$|$F) + PE classe válida &
Retorna \texttt{\{accessToken, refreshToken, perfil\}}; \textit{access token} verificável com \texttt{JWT\_SECRET} &
\texttt{findUser} mockado com \texttt{\{\_id, perfil, password\}}; credenciais correctas \\

\rowcolor{white}
TU-24 & RF-01, RF-02, RF-03 &
\texttt{POST /auth/login} &
Unidade &
PE — \texttt{user.id} sem \texttt{\_id} &
Retorna tokens; payload contém o \texttt{id} numérico e \texttt{perfil = Administrador} &
\texttt{findUser} mockado com \texttt{\{id, perfil, password\}} (sem campo \texttt{\_id}) \\

\midrule
% ==============================================================
% TI-06  —  hasProfile
% ==============================================================
\rowcolor{grouprow}
\multicolumn{7}{l}{%
  \footnotesize\textbf{TI-06 — \texttt{hasProfile(userPerfil, allowedProfiles)}}
  \quad Condição analisada: \texttt{!userPerfil || !allowedProfiles || !Array.isArray(allowedProfiles)}}\\
\midrule

\rowcolor{white}
TU-25 & RF-06 &
Rotas protegidas com \texttt{checkRole} &
Unidade &
CM\,—\,cond.\,(T) &
Devolve \texttt{false} &
Nenhuma \\

\rowcolor{rowalt}
TU-26 & RF-06 &
Rotas protegidas com \texttt{checkRole} &
Unidade &
CM\,—\,cond.\,(F$|$T) &
Devolve \texttt{false} &
\texttt{userPerfil} não nulo; \texttt{allowedProfiles} = \texttt{null} \\

\rowcolor{white}
TU-27 & RF-06 &
Rotas protegidas com \texttt{checkRole} &
Unidade &
CM\,—\,cond.\,(F$|$F$|$T) &
Devolve \texttt{false} &
\texttt{allowedProfiles} é \textit{string}, não \textit{array} \\

\rowcolor{rowalt}
TU-28 & RF-06 &
Rotas protegidas com \texttt{checkRole} &
Unidade &
CM (F$|$F$|$F) + PE classe válida-A &
Devolve \texttt{true} &
\texttt{userPerfil = Técnico}; \texttt{allowedProfiles = [Técnico, Responsável]} \\

\rowcolor{white}
TU-29 & RF-06 &
Rotas protegidas com \texttt{checkRole} &
Unidade &
CM (F$|$F$|$F) + PE classe válida-B &
Devolve \texttt{false} &
\texttt{userPerfil = Técnico}; \texttt{allowedProfiles = [Responsável, Administrador]} \\

\end{longtable}
\end{landscape}

% =============================================================
\section{Resumo de Cobertura}
% =============================================================

Resultados obtidos com \texttt{npx jest -{}-coverage} — todos os 29 casos \textbf{PASS}.

\vspace{0.4cm}

\begin{tabular}{@{}L{4.5cm}C{2.2cm}C{1.6cm}C{1.6cm}C{1.8cm}C{1.6cm}C{1.6cm}@{}}
  \toprule
  \textbf{Função} & \textbf{TU IDs} & \textbf{Casos} &
  \textbf{\% Stmts} & \textbf{\% Branch} & \textbf{\% Funcs} & \textbf{\% Lines} \\
  \midrule
  \rowcolor{rowalt}
  \texttt{generateAccessToken}  & TU-01–04 & 4  & 100\,\% & 100\,\% & 100\,\% & 100\,\% \\
  \texttt{generateRefreshToken} & TU-05–09 & 5  & 100\,\% & 100\,\% & 100\,\% & 100\,\% \\
  \rowcolor{rowalt}
  \texttt{verifyToken}          & TU-10–14 & 5  & 100\,\% & 100\,\% & 100\,\% & 100\,\% \\
  \texttt{renewAccessToken}     & TU-15–17 & 3  & 100\,\% & 100\,\% & 100\,\% & 100\,\% \\
  \rowcolor{rowalt}
  \texttt{login}                & TU-18–24 & 7  & 100\,\% & 100\,\% & 100\,\% & 100\,\% \\
  \texttt{hasProfile}           & TU-25–29 & 5  & 100\,\% & 100\,\% & 100\,\% & 100\,\% \\
  \midrule
  \rowcolor{grouprow}
  \textbf{Total} & TU-01–29 & \textbf{29} &
  \textbf{100\,\%} & \textbf{100\,\%} & \textbf{100\,\%} & \textbf{100\,\%} \\
  \bottomrule
\end{tabular}

\end{document}
